\documentclass[10pt,a4paper]{amsart}
\usepackage[margin=19mm]{geometry}
\usepackage{amsmath,amssymb,mathtools}
\usepackage[scr=rsfs,cal=euler]{mathalfa}
\usepackage{xfrac}
\usepackage[unicode,hidelinks,bookmarks=false]{hyperref}
\newcommand{\ie}{i.e.\ }
\newcommand{\cf}{cf.\ }
\newcommand{\resp}{respectively\ }
\newcommand{\Subsection}{\subsection}
\providecommand{\nobreakdash}{\nobreak\hbox{-}}
\setlength{\emergencystretch}{2em}
\multlinegap0pt
\usepackage{amssymb}
\usepackage[shortlabels]{enumitem}
\usepackage{xcolor}
\usepackage{mathtools}
\newtheorem{thm}{Theorem}
\newtheorem{prop}{Proposition}
\usepackage{cleveref}
\crefname{thm}{Theorem}{Theorems}
\crefname{prop}{Proposition}{Propositions}
\newcommand{\Div}{\operatorname{Div}}
\newcommand{\Frac}{\operatorname{Frac}}
\newcommand{\Ker}{\operatorname{Ker}}
\newcommand{\N}{\mathbb{N}}
\newcommand{\cO}{\mathcal{O}}
\newcommand{\dlog}{\operatorname{dlog}}
\newcommand{\hra}{\hookrightarrow}
\newcommand{\ra}{\rightarrow} 
\newcommand{\red}{{\operatorname{red}}}
\newcommand{\ssm}{\mathfrak{m}}
\renewcommand{\tilde}{\widetilde}
\newcommand{\ul}{\underline}
\title{Structure theorem for log de Rham-Witt sheaves with vanishing}
\author{Fei Ren}
\date{Source: arXiv:2501.06944v3 (CC BY 4.0)}
\begin{document}
\maketitle

\noindent\emph{Source and licence.} Structure theorem for log de Rham-Witt sheaves with vanishing. Version arXiv:2501.06944v3 (CC BY 4.0), distributed by arXiv under the Creative Commons Attribution 4.0 International licence, http://creativecommons.org/licenses/by/4.0/, as recorded in the arXiv licence metadata for this exact version and shown on the abstract page https://arxiv.org/abs/2501.06944v3. Full licence notice: \texttt{licenses/arxiv-2501.06944-CC-BY-4.0.txt}.\par\smallskip

\noindent\textit{Editorial scope.} Counted unit: the article's thm thm1 (source0.tex lines 527--557). The article's complete original proof of it is delivered with it (source0.tex lines 593--597, one \texttt{proof} environment of 5 lines, which the article places away from the statement). No mathematics is written, repaired or re-derived: the statement and the proof are the article's own lines, in the article's own notation, and no retained line is substituted or re-flowed.  Retained uncounted as the source-local context the proof needs: lines 600--643.  Nothing was omitted from any retained span, so the package carries no omission record.  No retained line contains a citation, so the wrapper carries no bibliography and no bibliography entry.  No retained line contains a figure environment, an image-inclusion command or drawing code, so no image is delivered and the package declares no asset dependency.  Editorial formatting only, in addition: an AMS \texttt{amsart} preamble that restates the article's own declarations and macros used by the retained lines, namely the environments \texttt{thm}, \texttt{prop} on separate counters, together with the standard packages those lines need.  The retained environments print in this document as Theorem 1, Proposition 1 in order of appearance, whereas the article prints the same items with the numbering of its own full document.  The complete native source of the article as distributed in the arXiv e-print for this exact version is bundled unchanged as \texttt{sources/171520/source0.tex}.
\begin{prop}
\label{MainProp}
Let $(R,\ssm,k')$ be the henselization of a closed point $x$ in $X$, where $\ssm$ is the maximal ideal of $R$, and $k'$ is the residue field.
Let $T_1,\cdots , T_d$ be a system of regular parameters of $R$ such that 
$$A=\Div(T_1\cdots T_e), \quad
B=\Div(T_1^{r_1}\cdots   T_g^{r_g})$$
where $e,f,g$ are integers such that $0\le e\le f \le g\le d$  and $r_1,\cdots, r_g$ are nonnegative integers such that at least one $r_i\ge 1$, and
$$\text{$p\nmid r_j\ge 1$ for all $j\in [e+1,f]$,
\quad 
$p\mid r_j\ge 1$ for all $j\in [f+1,g]$.}$$
Put 
$r_{g+1}=\cdots =r_{d}=0$ and set
$$\ul r:=(r_1,\cdots, r_d)\in \N^d.$$
Set
$$G_{R}^{\ul r}:=
\Omega^q_R(\log A)(-B),\quad
G_{R,\log}^{\ul r}:=
\Omega^q_R(\log A)(-B)
\cap
\Omega^q_{\Frac R,\log}.$$
% (Clearly these two abelian groups also depend on $A$, but we omit this from the notation.)
Define
$$\tilde r_i=
\begin{cases}
    r_i+1& i\in[e+1,f];
    \\
    r_i& i\notin [e+1,f].
\end{cases}$$
Set $\tilde{\underline{r}}:=(\tilde r_1,\cdots,\tilde r_d)$.
Then 
\begin{align}
\label{MainPropeq0}
G_{R,\log}^{\ul r}=
\Big\{
\sum_j^{\text{finite}}\dlog x_{1,j}\wedge &\dlog x_{2,j}\wedge\cdots \wedge\dlog x_{q,j}
\Big|
\\
&x_{1,j}\in 
1+(T_1^{\tilde r_1}\cdots  T_g^{\tilde r_g})\cdot R,\,
x_{2,j},\cdots,x_{q,j} \in R[\frac{1}{T_1\cdots T_f}]^\times
\Big\}.
\nonumber
\end{align}
\end{prop}

\begin{thm}
\label{thm1}
Let $A, B\ge 0$ be Cartier divisors on $X$ with $A$ and $(A+B)_\red$ being SNCDs. 
Write $$B=B'+B''$$ 
where 
$B'_\red\le A$, $B''$ and $A$ share no irreducible components, and $B'_\red+B''_\red$ being a SNCD. 
Let $$B''=B''_0+pB''_1$$
be the $p$-divisibility decomposition of $B''$ of length 1. 
Denote by $j:U=X\setminus (A+B)\hra X$ the open immersion.
For $q\ge 1$, let
$$\Omega^q_X(\log A)(-B)_{\log}:=
j_*\Omega^q_{U,\log}\cap \Omega^q_X(\log A)(-B).
$$
Let
$$\dlog:j_*(\underbrace{\cO_{U}^\times \otimes \cdots \otimes \cO_{U}^\times}_{q\text{-times}})
\ra j_*\Omega^q_U,\quad 
a_1\otimes \cdots\otimes a_q\mapsto \frac{da_1\wedge \cdots\wedge da_q}{a_1\cdots a_q}$$
be the $\dlog$ map.
Then 
we have the following identity of Nisnevich subsheaves of $\Omega^q_X(\log (A+B))$
$$\Omega^q_X(\log A)(-B)_{\log}=
\dlog
\left(
\Ker(\cO_X^\times \ra \cO_{\tilde B}^\times)\otimes 
\underbrace{ j_{0*}\cO_{V}^\times\otimes\cdots \otimes  j_{0*}\cO_{V}^\times}_{(q-1)\text{ times}}
\right).$$
where 
$$\tilde B=B+B''_{0,\red},$$
and
$j_0:V:=X\setminus (A+B''_0)\hra X$ is the open immersion.
\end{thm}

\begin{proof}[{Proof of \Cref{thm1}}]
Since $X$ is of finite type over $k$, it suffices to check that the two subsheaves have the same Nisnevich stalk at each closed point of $X$.
Let $R$ be the henselization of a local ring at a closed point of $X$.
\Cref{thm1} now follows from \Cref{MainProp}.
\end{proof}

\end{document}
